\documentclass[a4,10pt]{aleph-notas}

% -- Paquetes adicionales
\usepackage{enumitem}
\usepackage{aleph-comandos}

% -- Datos
\institucion{Proyecto Alephsub0}
\asignatura{Matemáticas de la Información}
\tema{Capacidad de un canal con regenerador ternario}
\autor{Andrés Merino}
\fecha{Abril 2026}

\fuente{montserrat}
\logouno[4.5cm]{Logos/LogoAlephsub0-02}
\definecolor{colordef}{cmyk}{0.81,0.62,0.00,0.22}

% -- Comandos adicionales
\newtheorem*{prob}{Problema}
    \tcolorboxenvironment{prob}{%
        color=colordef,recuadrost,colback=colordef!7,drop fuzzy shadow
    }

\begin{document}


\encabezado

\begin{prob}
Consideremos un sistema de transmisión de datos con entrada
\(X\in\{1,0,-1\}\), donde las probabilidades de los símbolos de entrada son:
\[
    P(X=1)=\alpha,\qquad P(X=0)=1-2\alpha,\qquad P(X=-1)=\alpha,
\]
con \(0<\alpha\le \frac{1}{2}\). El regenerador mantiene los símbolos no borrados y reconstruye los borrados
\((X=0)\) en proporción a la frecuencia de \(1\) y \(-1\).

Determinar la capacidad del canal (en bits por símbolo).
\end{prob}

\begin{proof}
Dado que el regenerador mantiene los símbolos no borrados y reconstruye los borrados en proporción a la frecuencia, se tiene:
\begin{align*}
    P(Y=1\mid X=1)&=1,& P(Y=-1\mid X=1)&=0,\\
    P(Y=1\mid X=-1)&=0,& P(Y=-1\mid X=-1)&=1,\\
    P(Y=1\mid X=0)&=\frac{\alpha}{2\alpha}=\frac{1}{2},&
    P(Y=-1\mid X=0)&=\frac{\alpha}{2\alpha}=\frac{1}{2}.
\end{align*}
Equivalente en forma matricial (filas: \(X=1,0,-1\); columnas: \(Y=1,-1\)):
\[
    \begin{pmatrix}
    1 & 0 \\
    \frac{1}{2} & \frac{1}{2} \\
    0 & 1
    \end{pmatrix}.
\]

Por el segundo teorema de Shannon, la capacidad es
\[
    C=\max_{P_X} I(X;Y).
\]
Para calcular la información mutua usamos
\[
    I(X;Y)=H(Y)-H(Y\mid X).
\]

Primero obtenemos la distribución de salida. Por la ley de la probabilidad total,
\begin{align*}
    P(Y=1)
        &=P(Y=1\mid X=1)P(X=1)+P(Y=1\mid X=0)P(X=0)+P(Y=1\mid X=-1)P(X=-1)\\
        &=1\cdot \alpha+\frac{1}{2}(1-2\alpha)+0\cdot \alpha\\
        &=\alpha+\frac{1}{2}-\alpha\\
        &=\frac{1}{2}.
\end{align*}
De manera análoga,
\[
P(Y=-1)=\frac{1}{2}.
\]

Así, \(Y\) es uniforme en \(\{1,-1\}\), y por tanto
\[
    H(Y)=-\frac{1}{2}\log_2\left(\frac{1}{2}\right)-\frac{1}{2}\log_2\left(\frac{1}{2}\right)=1.
\]

Calculamos ahora la entropía condicional para cada valor de \(X\):
\begin{itemize}
    \item Si \(X=0\), entonces \(Y\) toma los valores \(1\) y \(-1\) con probabilidad \(1/2\) cada uno, luego
    \begin{align*}
    H(Y\mid X=0)
        &= - P(Y=1\mid X=0)\log_2\left( P(Y=1\mid X=0)\right) \\
        & \phantom{=}- P(Y=-1\mid X=0)\log_2\left( P(Y=-1\mid X=0) \right)\\
        &= - \frac{1}{2}\log_2\left(\frac{1}{2}\right) 
        - \frac{1}{2}\log_2\left(\frac{1}{2}\right) = 1.
    \end{align*}
    \item Si \(X=1\), entonces \(Y=1\) con probabilidad \(1\), luego no hay incertidumbre:
    \[
        H(Y\mid X=1)=0.
    \]
    \item Si \(X=-1\), entonces \(Y=-1\) con probabilidad \(1\), luego tampoco hay incertidumbre:
    \[
        H(Y\mid X=-1)=0.
    \]

\end{itemize}

Por la fórmula de entropía condicional,
\begin{align*}
    H(Y\mid X)
        &=P(X=1)H(Y\mid X=1)+P(X=0)H(Y\mid X=0)+P(X=-1)H(Y\mid X=-1)\\
        &=\alpha\cdot 0+(1-2\alpha)\cdot 1+\alpha\cdot 0\\
        &=1-2\alpha.
\end{align*}

Sustituyendo en la expresión de información mutua,
\begin{align*}
    I(X;Y)
        &=H(Y)-H(Y\mid X)\\
        &=1-(1-2\alpha)\\
        &=2\alpha.
\end{align*}

Con esto, el valor máximo de $I$ en el intervalo \(0<\alpha\le \frac{1}{2}\) se alcanza en $\alpha=\frac{1}{2}$, en ese caso,
\[
    I(X;Y)=2\cdot \frac{1}{2}=1.
\]
Por consiguiente, la capacidad del canal es
\[
    C=1\text{ bit por símbolo}.\qedhere
\]
\end{proof}


\end{document}